\documentclass[11pt,a4paper]{article}
\usepackage{fontspec,amsmath,amssymb,amsthm,microtype}
\setmainfont{lmroman10-regular.otf}[BoldFont=lmroman10-bold.otf,ItalicFont=lmroman10-italic.otf,BoldItalicFont=lmroman10-bolditalic.otf]
\usepackage[margin=25mm]{geometry}
\usepackage[unicode,hidelinks]{hyperref}
\hypersetup{pdfauthor={LangevinLean Project}}
\setlength{\emergencystretch}{3em}
\setlength{\parskip}{0.4em}
\setlength{\parindent}{0pt}
\allowdisplaybreaks[2]
\raggedbottom
\title{\texorpdfstring{Assumptions and verification of the ULA convergence results}{Assumptions and verification of the ULA convergence results}}
\author{LangevinLean Project}
\date{}
\begin{document}\maketitle
Draft mathematical exposition; these notes may contain mistakes.


The convergence statements assume the specified potential regularity, step range, and initial integrability. Deterministic contraction and local-defect estimates give the mean-square bound. The construction of independent Gaussian innovations identifies the numerical marginal laws, and taking the initial coupling infimum gives the Wasserstein bound.



Results stated with a supplied reference or a weak-equation uniqueness hypothesis retain those assumptions. An application must either provide them or use the subsequent theorems that prove the required reference properties. The convergence statements are therefore read together with their complete hypotheses.


\section{\texorpdfstring{Verification and mathematical hypotheses}{Verification and mathematical hypotheses}}


The convergence result separates deterministic geometry, probabilistic construction, and the final transport comparison. Each application must verify the hypotheses at the appropriate level: gradient regularity for contraction, target moments and stationarity for local error, and independent innovations for identifying the numerical laws. Conditional weak-equation statements retain their uniqueness hypotheses.


\end{document}
